\documentclass{teipel-thesis-xelatex}
\usepackage[utf8]{inputenc}
\usepackage[T1]{fontenc}
\usepackage{geometry}
\usepackage{graphicx}
\usepackage{float}
\usepackage{hyperref}
\usepackage[english,indonesian]{babel}

\title{Optimasi Pemrosesan Template LaTeX Menggunakan Pendekatan Berbasis AI}
\author{Academic Support Agent}
\date{\today}

\begin{document}
\maketitle

\begin{abstract}

\end{abstract}

\chapter{PENDAHULUAN}


\chapter{TINJAUAN PUSTAKA}


% Original content (commented out):
% %%%%%%%%%%%%%%%%%%%%%%%%%%%%%%%%%%%%%%%%%%%%%%%%%%%%%%%%%%%%%%%%%%%%%%
%% File: main.tex
%% Author: John Liaperdos
%% Last update: Sept 29, 2015
%% Description: Provides an example of a Diploma Thesis using the
%% teipel-thesis-xelatex class.
%%
%% -------------------------------------------------------------------
%% WARNING: If an index is to be typeset, file name is allowed to contain 
%% only one dot, e.g. "main.tex" is an acceptable name, while 
%% "main.file.tex" is not.
%% -------------------------------------------------------------------
%%
%%%%%%%%%%%%%%%%%%%%%%%%%%%%%%%%%%%%%%%%%%%%%%%%%%%%%%%%%%%%%%%%%%%%%%
%% CLASS OPTIONS:
%%%%%%%%%%%%%%%%%%%%%%%%%%%%%%%%%%%%%%%%%%%%%%%%%%%%%%%%%%%%%%%%%%%%%%
%%
%% 1. add/remove the "hyperref" option to enable/disable hyperlinks.
%% (NOTE: It might be useful to delete previously existing auxiliary
%% files before re-compiling after the removal of the "hyperref" 
%% option.)
%%
%% 2. add/remove the "printer" option to typeset a printer-friendly 
%% (grayscale)/color version of the thesis.
%%
%% 3. use the "watermark" option to indicate that this is not an actual
%% thesis.
%%
%% 4. use the "histinit" option to enable "historiated initials".
%% (If used, all chapter initials declared by the \InitialCharacter{}
%% macro are enlarged. If omitted, arguments of \InitialCharacter{}
%% are typeset as normal text.)
%%
%% 5. use the "arial/times/clearsans" option to typeset using the 
%% respective fonts.
%%
%% 6. use the "noindex" option to (hopefully) avoid compilation timeouts
%% when compiling online (disables index generation - note that "\indexGR"
%% invocations need not be removed when toggling this option).
%%
%% 7. use the "plain" option to disable tikz graphics in title page
%% and part/chapter headers (might help to avoid compilation timeouts).
%% Note that "plain" disables CD label and cover creation.
%%
%% 8. activate the "newlogo" option to use the new official Logo.
%%
%%%%%%%%%%%%%%%%%%%%%%%%%%%%%%%%%%%%%%%%%%%%%%%%%%%%%%%%%%%%%%%%%%%%%%
%%
\documentclass[clearsans,hyperref,watermark,%
		histinit,noindex,newlogo,plain]{teipel-thesis-xelatex}
%%
%%%%%%%%%%%%%%%%%%%%%%%%%%%%%%%%%%%%%%%%%%%%%%%%%%%%%%%%%%%%%%%%%%%%%%
%%
%%
%%%%%%%%%%%%%%%%%%%%%%%%%%%%%%%%%%%%%%%%%%%%%%%%%%%%%%%%%%%%%%%%%%%%%%
%% THESIS INFO: 
%%%%%%%%%%%%%%%%%%%%%%%%%%%%%%%%%%%%%%%%%%%%%%%%%%%%%%%%%%%%%%%%%%%%%%
%%
%% ΤΙΤΛΟΣ ΠΤΥΧΙΑΚΗΣ ΕΡΓΑΣΙΑΣ 
%%
%% Για εξαναγκασμένες αλλαγές γραμμής χρησιμοποιήστε "\\".
%% Αν οι αλλαγές γραμμής πρέπει να είναι διαφορετικές στο εξώφυλλο σε σχέση 
%% με το εσώφυλλο (σελ. 3), επαναλάβετε τον τίτλο του εξωφύλλου με τις 
%% επιθυμητές αλλαγές γραμμής ως προαιρετικό όρισμα της εντολής \title.
%%
%% Παραδείγματα:
%% 1. Όμοιος τίτλος σε εξώφυλλο και εσώφυλλο, με αυτόματες αλλαγές γραμμής:
%%	    \title{Πρότυπο Σύστημα Ομότιμων Κόμβων Βασισμένο σε Σχήματα RDF}
%% 2. Όμοιος τίτλος σε εξώφυλλο και εσώφυλλο, με αλλαγή γραμμής μετά τη λέξη
%% "Σύστημα":
%%	    \title{Πρότυπο Σύστημα \\ Ομότιμων Κόμβων Βασισμένο σε Σχήματα RDF}
%% 3. Διαφορετικές αλλαγές γραμμής σε εξώφυλλο και εσώφυλλο. Στο εξώφυλλο 
%% έχουμε αλλαγή γραμμής μετά τη λέξη "Σύστημα", ενώ στο εσώφυλλο η αλλαγή
%% γραμμής ακολουθεί τη λέξη "Ομότιμων":
%%	    \title[Πρότυπο Σύστημα \\ Ομότιμων Κόμβων Βασισμένο σε Σχήματα RDF]% (προαιρετικό όρισμα)
%%            {Πρότυπο Σύστημα Ομότιμων \\ Κόμβων Βασισμένο σε Σχήματα RDF}% (υποχρεωτικό όρισμα)
%%
	\title{Πρότυπο Σύστημα Ομότιμων Κόμβων \\Βασισμένο σε Σχήματα RDF} 
%%
%% -------------------------------------------------------------------
%% ΥΠΟΤΙΤΛΟΣ ΠΤΥΧΙΑΚΗΣ ΕΡΓΑΣΙΑΣ (προαιρετικός)
%%
%% Αν δεν υπάρχει υπότιτλος, τοποθετήστε τον χαρακτήρα του σχολίου "%"
%% πριν από την εντολή \subtitle, ή αφήστε κενό το όρισμα της εντολής.
%%
%% Παράδειγμα:
%%	\subtitle{Μελέτη και υλοποίηση}
	\subtitle{Μελέτη και υλοποίηση}
%%
%% -------------------------------------------------------------------
%% ΛΕΞΕΙΣ ΚΛΕΙΔΙΑ, ΣΤΑ ΕΛΛΗΝΙΚΑ
%%
%% Παράδειγμα: \keywordsGR{Σύστημα ομότιμων κόμβων, Σύστημα ομότιμων κόμβων βασισμένο σε
%% σχήματα, Σημασιολογικός Ιστός, RDF/S, RQL, Jxta}
	\keywordsGR{Σύστημα ομότιμων κόμβων, Σύστημα ομότιμων κόμβων βασισμένο %
    	σε σχήματα, Σημασιολογικός Ιστός, RDF/S, RQL, Jxta}
%%
%% -------------------------------------------------------------------
%% ΛΕΞΕΙΣ ΚΛΕΙΔΙΑ, ΣΤΑ ΑΓΓΛΙΚΑ
%%
%% Παράδειγμα: \keywordsEN{Peer-to-peer, Schema-based peer-to-peer, Semantic Web, 
%% RDF/S, RQL, Jxta}
	\keywordsEN{Peer-to-peer, Schema-based peer-to-peer, Semantic Web, % 
    	RDF/S, RQL, Jxta}
%%
%% -------------------------------------------------------------------
%% ΘΕΜΑΤΙΚΗ ΠΕΡΙΟΧΗ
%%
%% (εμφανίζεται μόνο στις ιδιότητες του παραγόμενου αρχείου PDF)
%% Αν δεν είναι επιθυμητή, τοποθετήστε τον χαρακτήρα του σχολίου "%"
%% πριν από την εντολή \subjectArea, ή αφήστε κενό το όρισμα της εντολής.
%%
%% Παράδειγμα: \subjectArea{Σχήματα RDF/S}
	\subjectArea{Σχήματα RDF/S}
%%
%% -------------------------------------------------------------------
%% ΤΟΥ/ΤΗΣ/ΤΩΝ
%%
%% "του" ή "της" ή "των", ανάλογα με το φύλο/αριθμό του σπουδαστή ή 
%% των σπουδαστών
%%
%% Παράδειγμα:
%%	\toutis{του}
	\toutis{του}
%%
%% -------------------------------------------------------------------
%% ΟΝΟΜΑΤΕΠΩΝΥΜΟ ΣΠΟΥΔΑΣΤΗ (ΚΕΦΑΛΑΙΑ, ΓΕΝΙΚΗ ΠΤΩΣΗ)
%%
%% Για περισσότερους του ενός σπουδαστές, διαχωρίστε με ",".
%%
%% Παράδειγμα:
%%	\authorNameCapital{ΚΩΝΣΤΑΝΤΙΝΟΥ Δ. ΔΗΜΗΤΡΙΟΥ, ΓΕΩΡΓΙΟΥ Π. ΠΑΝΑΓΑΚΗ}
	\authorNameCapital{ΚΩΝΣΤΑΝΤΙΝΟΥ Δ. ΔΗΜΗΤΡΙΟΥ}
%%
%% -------------------------------------------------------------------
%% ΟΝΟΜΑΤΕΠΩΝΥΜΟ ΣΠΟΥΔΑΣΤΗ (ΠΕΖΑ, ΟΝΟΜΑΣΤΙΚΗ ΠΤΩΣΗ)
%% 
%% Για περισσότερους του ενός σπουδαστές, διαχωρίστε με ",".
%%
%% Παράδειγμα:
%%	\authorName{Κωνσταντίνος Δημητρίου, Γεώργιος Παναγάκης}
	\authorName{Κωνσταντίνος Δημητρίου}
%%
%% -------------------------------------------------------------------
%% ΟΝΟΜΑΤΕΠΩΝΥΜΟ ΕΠΙΒΛΕΠΟΝΤΑ ΚΑΘΗΓΗΤΗ
%% 
	\supervisor{Γιάννης Λιαπέρδος}
%%
%% -------------------------------------------------------------------
%% ΤΙΤΛΟΣ ΕΠΙΒΛΕΠΟΝΤΑ ΚΑΘΗΓΗΤΗ
	\supervisorTitle{Καθηγητής Εφαρμογών}
%%
%% -------------------------------------------------------------------
%% ΕΠΙΒΛΕΠΩΝ/ΕΠΙΒΛΕΠΟΥΣΑ
%%
%% "Επιβλέπων" ή "Επιβλέπουσα", ανάλογα με το φύλο του 
%% Επιβλέποντα Καθηγητή
	\supervisorMaleFemale{Επιβλέπων}
%%
%% -------------------------------------------------------------------
%% ΤΟΠΟΣ/ΜΗΝΑΣ/ΕΤΟΣ ΕΚΔΟΣΗΣ
%%
	\thesisPlaceDate{Σπάρτη, Νοέμβριος 2014}
%%
%% -------------------------------------------------------------------
%% ΤΟΠΟΣ/ΜΗΝΑΣ/ΕΤΟΣ ΣΥΓΓΡΑΦΗΣ (Εμφανίζεται στη σελίδα των ευχαριστιών,
%% αν υπάρχει).
%%
	\ackPlaceDate{Κέρκυρα, Μάιος 2014}
%%
%% -------------------------------------------------------------------
%% ΗΜΕΡΟΜΗΝΙΑ ΕΞΕΤΑΣΗΣ
%%
	\examinationDate{22α Νοεμβρίου 2014}
%%
%% -------------------------------------------------------------------
%% ΗΜΕΡΟΜΗΝΙΑ ΔΗΛΩΣΗΣ ΠΕΡΙ ΜΗ ΛΟΓΟΚΛΟΠΗΣ
%%
	\declarationDate{22 Σεπτεμβρίου 2014}
%%
%% -------------------------------------------------------------------
%% ΕΤΟΣ COPYRIGHT
%%
	\copyrightYear{2014}
%%
%% -------------------------------------------------------------------
%% ΟΝΟΜΑΤΕΠΩΝΥΜΟ 1ου ΕΞΕΤΑΣΤΗ
%%
	\firstExaminer{Γρηγόρης Καραγιώργος}
%%
%% -------------------------------------------------------------------
%% ΤΙΤΛΟΣ 1ου ΕΞΕΤΑΣΤΗ
%%
	\firstExaminerTitle{Επίκουρος Καθηγητής}
%%
%% -------------------------------------------------------------------
%% ΟΝΟΜΑΤΕΠΩΝΥΜΟ 2ου ΕΞΕΤΑΣΤΗ
%%
	\secondExaminer{Γεώργιος Γεωργίου}
%%
%% -------------------------------------------------------------------
%% ΤΙΤΛΟΣ 2ου ΕΞΕΤΑΣΤΗ
%%
	\secondExaminerTitle{Επιστ. Συνεργάτης}
%%
%%
%%%%%%%%%%%%%%%%%%%%%%%%%%%%%%%%%%%%%%%%%%%%%%%%%%%%%%%%%%%%%%%%%%%%%%
%% THESIS COLORS: 
%%%%%%%%%%%%%%%%%%%%%%%%%%%%%%%%%%%%%%%%%%%%%%%%%%%%%%%%%%%%%%%%%%%%%%
%%
%% Χρώμα εξωφύλλου - κεφαλαίων
	\chaptercolor{gray!50!brown}
%%
%% Χρώμα παραρτημάτων
	\appendixcolor{brown!60!orange}
%%
%% Χρώμα υπερσυνδέσμων (αν έχει ενεργοποιηθεί η επιλογή "hyperref")
    \hyperlinkcolor{blue}
%%
%% Χρώμα τίτλου εργασίας στο εξώφυλλο (αν δεν έχει ενεργοποιηθεί 
%% η επιλογή "plain")
    \titlecolor{white}
%%
%% Χρώμα υποβάθρου (φόντου) τίτλου εργασίας στο εξώφυλλο (αν δεν έχει 
%% ενεργοποιηθεί η επιλογή "plain")
    \titlebackgroundcolor{gray!60!brown}  
%%
%%
%%%%%%%%%%%%%%%%%%%%%%%%%%%%%%%%%%%%%%%%%%%%%%%%%%%%%%%%%%%%%%%%%%%%%%
%% COVER PAGE IMAGE: 
%%%%%%%%%%%%%%%%%%%%%%%%%%%%%%%%%%%%%%%%%%%%%%%%%%%%%%%%%%%%%%%%%%%%%%
%%
%% Εικόνα εξωφύλλου (προαιρετική)
%% Στην περίπτωση κατά την οποία δεν είναι επιθυμητή η εισαγωγή εικόνας στο εξώφυλλο,
%% διαγράψτε την εντολή \coverpageimage, ή μετατρέψτε την σε σχόλιο (με "%")
%%
%% Σύνταξη:
%%          \coverpageimage{συντελεστής μεγέθυνσης}{όνομα αρχείου εικόνας [πλήρης διαδρομή]}
%%      ή
%%          \coverpageimage[tikz]{συντελεστής μεγέθυνσης}{εντολές TikZ}
%%          (στις εντολές μπορούν να περιλαμβάνονται και δηλώσεις \usetikzlibrary, κ.λπ.)
%%      
%% Παραδείγματα:
%%      - Χρήση εικόνας από το αρχείο "figures/rdf.png" με συντελεστή μεγέθυνσης 0.8:
%%          \coverpageimage{0.8}{figures/rdf.png}
%%      - Χρήση εικόνας TikZ με συντελεστή μεγέθυνσης 0.5:
%%          \coverpageimage[tikz]{0.5}{
%%              \draw[thick, magenta] \foreach \x in {18,90,...,306} {
%%                  (\x:4) node{} -- (\x+72:4)
%%                  (\x:4) -- (\x:3) node{}
%%                  (\x:3) -- (\x+15:2) node{}
%%                  (\x:3) -- (\x-15:2) node{}
%%                  (\x+15:2) -- (\x+144-15:2)
%%                  (\x-15:2) -- (\x+144+15:2)
%%              };
%%          }
%%
     \coverpageimage[tikz]{0.5}{
        \draw[thick, gray] \foreach \x in {18,90,...,306} {
            (\x:4) node{} -- (\x+72:4)
            (\x:4) -- (\x:3) node{}
            (\x:3) -- (\x+15:2) node{}
            (\x:3) -- (\x-15:2) node{}
            (\x+15:2) -- (\x+144-15:2)
            (\x-15:2) -- (\x+144+15:2)
        };
    }
%%
%%%%%%%%%%%%%%%%%%%%%%%%%%%%%%%%%%%%%%%%%%%%%%%%%%%%%%%%%%%%%%%%%%%%%%

\begin{document}

\maketitle

\beginfrontmatter
	
% Περίληψη
	\include{abstract}
% Αφιέρωση
	\thesisDedication{στους γονείς μου}
% Ευχαριστίες
	\include{acknowledgements}
% Πίνακας Περιεχομένων
	\tableofcontents
% Κατάλογος Σχημάτων
	\listoffigures
% Κατάλογος Εικόνων
	\listofillustrations
% Κατάλογος Πινάκων
	\listoftables
% Πρόλογος
	\include{preface}
	
\beginmainmatter

%%%%%%%%%%%%%%%%%%%%%%%%%%%%%%%%%%%%%%%%%%%%%%%%%%%%%%%%%%%%%%%%%%%%%%
%% INCLUDE YOUR CHAPTERS/SECTIONS HERE:
%%%%%%%%%%%%%%%%%%%%%%%%%%%%%%%%%%%%%%%%%%%%%%%%%%%%%%%%%%%%%%%%%%%%%%
%%
% Εισαγωγή
	\include{intro}
% Μέρη/Κεφάλαια
	\part{Θεωρητικό Μέρος}
		\include{chap2}
		\include{chap3}
	\part{Πρακτικό Μέρος} 
		\include{chap4}
		\include{chap5}
		\include{chap6}
		\include{chap7}
		\include{chap8}
	\part{Επίλογος}
		\include{chap9}
% Παραρτήματα
	\appendices
		\include{appA}
		\include{appB}
		\include{appC}
		\include{appD}
		\include{appE} 
% Βιβλιογραφία - Αναφορές
	\bibliography{references}
% Συντομογραφίες - Αρκτικόλεξα - Ακρωνύμια
	\includeabbreviations{abbreviations}
% Γλωσσάριο
	\includeglossary{glossary}
%%%%%%%%%%%%%%%%%%%%%%%%%%%%%%%%%%%%%%%%%%%%%%%%%%%%%%%%%%%%%%%%%%%%%%
% Ευρετήριο Όρων
	\printindices
%
%%%%%%%%%%%%%%%%%%%%%%%%%%%%%%%%%%%%%%%%%%%%%%%%%%%%%%%%%%%%%%%%%%%%%%
%%
%%
%%%%%%%%%%%%%%%%%%%%%%%%%%%%%%%%%%%%%%%%%%%%%%%%%%%%%%%%%%%%%%%%%%%%%%
%% CREATE CD COVERS AND LABELS:
%%%%%%%%%%%%%%%%%%%%%%%%%%%%%%%%%%%%%%%%%%%%%%%%%%%%%%%%%%%%%%%%%%%%%%
%%
%% Δημιουργία ετικετών CD:
%%
	\definecdlabeloffsets{0}{-0.65}{0}{0.55} 
%% Παράμετροι εντολής:
%   upper label x offset [cm] (default=0) /  upper label y offset [cm] (default=0) /  
%   lower label x offset [cm] (default=0) /  lower  label y offset [cm] (default=0) 
%   -- For Q-Connect KF01579 labels use the following offset values: {0}{-0.65}{0}{0.55}
%%
	\createcdlabel{Πρότυπο Σύστημα Ομότιμων \\Κόμβων Βασισμένο σε Σχήματα RDF}%
	    {Κωνσταντίνος Δ. Δημητρίου}{ΟΚΤΩΒΡΙΟΣ}{2014}{8} 
%% Παράμετροι εντολής:	
%   τίτλος πτυχιακής / όνομα συγγραφέα / μήνας / έτος / εύρος περιοχής τίτλου 
%   σε cm (προτεινόμενη τιμή: 8)
%% ΣΗΜΕΙΩΣΗ: Για εξαναγκασμένες αλλαγές γραμμής χρησιμοποιήστε "\\".
%%
%%
%% Δημιουργία εξωφύλλου θήκης CD:
%%
	\createcdcover{Πρότυπο Σύστημα Ομότιμων \\Κόμβων Βασισμένο σε Σχήματα RDF}%
	    {Κωνσταντίνος Δ. Δημητρίου}{ΟΚΤΩΒΡΙΟΣ}{2014}{10} 
%% Παράμετροι εντολής:		
%   τίτλος πτυχιακής / όνομα συγγραφέα / μήνας / έτος / εύρος περιοχής τίτλου 
%   σε cm (προτεινόμενη τιμή: 10) 
%% ΣΗΜΕΙΩΣΗ: Για εξαναγκασμένες αλλαγές γραμμής χρησιμοποιήστε "\\".
%%
%
%%
\end{document}
%%
%%%%%%%%%%%%%%%%%%%%%%%%%%%%%%%%%%%%%%%%%%%%%%%%%%%%%%%%%%%%%%%%%%%%%%

\end{document}